\documentclass[11pt,a4paper]{article}
\usepackage[utf8]{inputenc}
\usepackage[T1]{fontenc}
\usepackage[spanish,es-minimal,es-noshorthands]{babel}
\usepackage{lmodern}
\usepackage[sfdefault,lining]{FiraSans}
\usepackage[a4paper,top=2.3cm,bottom=2.2cm,left=1.8cm,right=1.8cm,headheight=1.3cm,headsep=0.4cm,footskip=1cm]{geometry}
\usepackage{amsmath,amssymb}
\allowdisplaybreaks
\usepackage{icomma}
\usepackage{graphicx}
\usepackage[table]{xcolor}
\usepackage{tikz}
\usepackage{fontawesome5}
\usepackage{ccicons}
\usepackage{fancyhdr}
\usepackage{lastpage}
\usepackage{tasks}
\usepackage{needspace}
\usepackage[most]{tcolorbox}
\usepackage{array}
\usepackage[hidelinks]{hyperref}
\AtBeginDocument{\hypersetup{pdftitle={Soluciones de la ficha de la unidad 2: Sistemas de ecuaciones lineales. 2º Bach.},pdfsubject={Matemáticas II 2º Bach.},pdfkeywords={sistemas de ecuaciones lineales, método de Gauss, método matricial, regla de Cramer, teorema de Rouché-Frobenius, sistemas homogéneos, discusión de sistemas, parámetros, problemas, PAU}}}

\definecolor{navy}{HTML}{1B2748}
\definecolor{accent}{HTML}{E8603A}
\definecolor{teal}{HTML}{1F9E90}
\definecolor{soft}{HTML}{EEF1F8}
\definecolor{muted}{HTML}{5D6685}

\setlength{\parindent}{0pt}
\setlength{\parskip}{3pt}
\settasks{label=\alph*), label-format=\bfseries\color{navy}, label-width=1.3em, item-indent=1.8em, column-sep=0.6em, after-item-skip=5pt, before-skip=3pt, after-skip=4pt}

% ---------- capçalera i peu ----------
\newcommand{\llicencia}{\href{https://creativecommons.org/licenses/by-nc-sa/4.0/deed.es}{\ccbyncsa}\ \ Licencia Creative Commons Atribución-NoComercial-CompartirIgual 4.0 Internacional (CC BY-NC-SA 4.0)}
\pagestyle{fancy}
\fancyhf{}
\fancyhead[L]{\raisebox{-0.35\height}{\includegraphics[height=0.95cm]{logo-gva.png}}}
\fancyhead[C]{\small\color{muted}Soluciones de la ficha · Unidad 2: Sistemas de ecuaciones lineales · 2\textsuperscript{o} Bach.}
\fancyhead[R]{\raisebox{-0.35\height}{\includegraphics[height=0.95cm]{logo-vermellar.png}}}
\fancyfoot[L]{\scriptsize\color{muted}\llicencia}
\fancyfoot[R]{\small\color{muted}\thepage}
\renewcommand{\headrulewidth}{0.4pt}
\renewcommand{\headrule}{\color{muted!40}\hrule height \headrulewidth}
\fancypagestyle{primera}{\fancyhf{}\renewcommand{\headrulewidth}{0pt}
  \fancyfoot[L]{\scriptsize\color{muted}\llicencia}
  \fancyfoot[R]{\small\color{muted}\thepage}}

% ---------- nivells i exercicis ----------
\newcommand{\estrelles}[1]{{\color{accent}\ifcase#1\or\faStar\faStar[regular]\faStar[regular]\or\faStar\faStar\faStar[regular]\or\faStar\faStar\faStar\fi}}
\newcounter{exer}
\newcommand{\ex}[1]{\par\needspace{4\baselineskip}\medskip\refstepcounter{exer}\noindent{\color{navy}\bfseries\theexer.}\ \ }

\newcommand{\seccio}[2]{\par\needspace{10\baselineskip}\bigskip
  \begin{tcolorbox}[colback=navy,colframe=navy,coltext=white,arc=2mm,boxrule=0pt,left=3mm,right=3mm,top=1.2mm,bottom=1.2mm]
  {\large\bfseries #1\quad #2}
  \end{tcolorbox}\vspace{-2pt}}


% quadre de recordatori (només a la fitxa d'exercicis)
\newcommand{\recorda}[1]{\par\smallskip\begin{tcolorbox}[colback=accent!7,colframe=accent!55,arc=2mm,boxrule=0.5pt,left=2.5mm,right=2.5mm,top=1mm,bottom=1mm]\small{\color{accent}\faLightbulb[regular]}\ \textbf{Recuerda:}\ #1\end{tcolorbox}}
% etiqueta dels exercicis de la PAU d'altres comunitats
\newcommand{\pau}[1]{{\color{teal}\small\bfseries\faUniversity\ PAU · #1.}\ }
\newcommand{\res}[1]{\boxed{#1}}\newcommand{\rest}[1]{\fbox{#1}}

% =====================================================================
% Capçalera FentMates de les fitxes: un color per a cada curs.
% Els colors són els de la pàgina inicial (LEVELS a assets/js/materials.js);
% el logo ve de _eines/logo/pdf/ (python3 _eines/logo/genera.py).
%
% Ús, al final del preàmbul de la unitat:
%   % =====================================================================
% Capçalera FentMates de les fitxes: un color per a cada curs.
% Els colors són els de la pàgina inicial (LEVELS a assets/js/materials.js);
% el logo ve de _eines/logo/pdf/ (python3 _eines/logo/genera.py).
%
% Ús, al final del preàmbul de la unitat:
%   % =====================================================================
% Capçalera FentMates de les fitxes: un color per a cada curs.
% Els colors són els de la pàgina inicial (LEVELS a assets/js/materials.js);
% el logo ve de _eines/logo/pdf/ (python3 _eines/logo/genera.py).
%
% Ús, al final del preàmbul de la unitat:
%   \input{../comu/fentmates.tex}
%   \fmcurs{1eso}
% i, al començament del document:
%   \capcalera{Matemàtiques}{1r ESO}{Fitxa d'exercicis · Unitat 2}{Divisibilitat}
% Opcional abans de \input: \def\fmneed{14} (espai mínim abans de cada secció).
% =====================================================================
\makeatletter
\def\fm@logodir{../../logo/pdf}
\providecommand{\fmneed}{10}
\def\fm@cursnom{}

% clau del curs, color i color del text damunt del color
\newcommand{\fm@curs}[3]{\@namedef{fm@color@#1}{#2}\@namedef{fm@text@#1}{#3}}
\fm@curs{1eso}{E6B84A}{navy}
\fm@curs{2eso}{3B82F6}{white}
\fm@curs{3eso}{FF7A50}{white}
\fm@curs{4eso-a}{A78BFA}{white}
\fm@curs{4eso-b}{6D28D9}{white}
\fm@curs{1batx-mi}{9B1C1C}{white}
\fm@curs{1batx-ccss}{EF6B6B}{white}
\fm@curs{2batx-mii}{1E6B47}{white}
\fm@curs{2batx-ccss}{2BB3A3}{white}

\newcommand{\fmcurs}[1]{%
  \@ifundefined{fm@color@#1}{\PackageError{fentmates}{Curs desconegut: #1}{}}{}%
  \edef\fm@tmp{\noexpand\definecolor{curs}{HTML}{\@nameuse{fm@color@#1}}%
               \noexpand\colorlet{curstext}{\@nameuse{fm@text@#1}}}\fm@tmp
  \colorlet{cursfosc}{curs!72!black}%
  \def\fm@logo{\fm@logodir/vertical-#1.pdf}%
  \def\fm@marca{\fm@logodir/marca-#1.pdf}}

% ---------- capçalera de la primera pàgina ----------
% #1 matèria (petit), #2 curs, #3 tipus de document i unitat, #4 títol
\newcommand{\capcalera}[4]{%
  \gdef\fm@cursnom{#2}%
  \par\noindent
  \begin{tcolorbox}[colback=curs,colframe=curs,boxrule=0pt,arc=2.2mm,width=3.5cm,height=2.15cm,
      valign=center,left=3mm,right=2mm,top=1mm,bottom=1mm,nobeforeafter,box align=center]
    \color{curstext}{\footnotesize #1}\par\vspace{1pt}{\fontsize{19}{21}\selectfont\bfseries #2\par}%
  \end{tcolorbox}\hspace{5mm}%
  \begin{minipage}[c]{\dimexpr\linewidth-3.5cm-5mm-2.9cm\relax}
    {\color{cursfosc}\small\bfseries\MakeUppercase{#3}}\par\vspace{2pt}
    {\color{navy}\fontsize{20}{23}\selectfont\bfseries #4\par}
  \end{minipage}\hfill
  \begin{minipage}[c]{2.6cm}\raggedleft\includegraphics[width=2.6cm]{\fm@logo}\end{minipage}\par
  \vspace{2pt}}

% ---------- barres de secció del color del curs ----------
\renewcommand{\seccio}[2]{\par\needspace{\fmneed\baselineskip}\bigskip
  \begin{tcolorbox}[colback=curs,colframe=curs,coltext=curstext,arc=2mm,boxrule=0pt,left=3mm,right=3mm,top=1.2mm,bottom=1.2mm]
  {\large\bfseries #1\quad #2}
  \end{tcolorbox}\vspace{-2pt}}

% ---------- estrelles de dificultat del color del curs (a Batxillerat no se'n posen) ----------
\@ifundefined{estrelles}{}{%
  \renewcommand{\estrelles}[1]{{\color{curs}\ifcase#1\or\faStar\faStar[regular]\faStar[regular]\or\faStar\faStar\faStar[regular]\or\faStar\faStar\faStar\fi}}}

% ---------- capçalera i peu de les pàgines ----------
% primera pàgina: logos institucionals al peu
\fancypagestyle{primera}{\fancyhf{}\renewcommand{\headrulewidth}{0pt}
  \fancyfoot[L]{\raisebox{-0.32\height}{\includegraphics[height=0.78cm]{logo-gva.png}}\hspace{5mm}%
    \raisebox{-0.32\height}{\includegraphics[height=0.78cm]{logo-vermellar.png}}}
  \fancyfoot[R]{\scriptsize\color{muted}\href{https://creativecommons.org/licenses/by-nc-sa/4.0/deed.ca}{\ccbyncsa}\ \ CC BY-NC-SA 4.0 · fentmates.org\qquad\small\thepage}}
% la resta: el símbol a l'esquerra i el curs a la dreta, amb el títol al centre com abans
\fancyhead[L]{\raisebox{-0.3\height}{\includegraphics[height=0.62cm]{\fm@marca}}}
\fancyhead[R]{\small\bfseries\color{cursfosc}\fm@cursnom}
\renewcommand{\headrule}{\color{curs!45}\hrule height \headrulewidth}
\makeatother

%   \fmcurs{1eso}
% i, al començament del document:
%   \capcalera{Matemàtiques}{1r ESO}{Fitxa d'exercicis · Unitat 2}{Divisibilitat}
% Opcional abans de \input: \def\fmneed{14} (espai mínim abans de cada secció).
% =====================================================================
\makeatletter
\def\fm@logodir{../../logo/pdf}
\providecommand{\fmneed}{10}
\def\fm@cursnom{}

% clau del curs, color i color del text damunt del color
\newcommand{\fm@curs}[3]{\@namedef{fm@color@#1}{#2}\@namedef{fm@text@#1}{#3}}
\fm@curs{1eso}{E6B84A}{navy}
\fm@curs{2eso}{3B82F6}{white}
\fm@curs{3eso}{FF7A50}{white}
\fm@curs{4eso-a}{A78BFA}{white}
\fm@curs{4eso-b}{6D28D9}{white}
\fm@curs{1batx-mi}{9B1C1C}{white}
\fm@curs{1batx-ccss}{EF6B6B}{white}
\fm@curs{2batx-mii}{1E6B47}{white}
\fm@curs{2batx-ccss}{2BB3A3}{white}

\newcommand{\fmcurs}[1]{%
  \@ifundefined{fm@color@#1}{\PackageError{fentmates}{Curs desconegut: #1}{}}{}%
  \edef\fm@tmp{\noexpand\definecolor{curs}{HTML}{\@nameuse{fm@color@#1}}%
               \noexpand\colorlet{curstext}{\@nameuse{fm@text@#1}}}\fm@tmp
  \colorlet{cursfosc}{curs!72!black}%
  \def\fm@logo{\fm@logodir/vertical-#1.pdf}%
  \def\fm@marca{\fm@logodir/marca-#1.pdf}}

% ---------- capçalera de la primera pàgina ----------
% #1 matèria (petit), #2 curs, #3 tipus de document i unitat, #4 títol
\newcommand{\capcalera}[4]{%
  \gdef\fm@cursnom{#2}%
  \par\noindent
  \begin{tcolorbox}[colback=curs,colframe=curs,boxrule=0pt,arc=2.2mm,width=3.5cm,height=2.15cm,
      valign=center,left=3mm,right=2mm,top=1mm,bottom=1mm,nobeforeafter,box align=center]
    \color{curstext}{\footnotesize #1}\par\vspace{1pt}{\fontsize{19}{21}\selectfont\bfseries #2\par}%
  \end{tcolorbox}\hspace{5mm}%
  \begin{minipage}[c]{\dimexpr\linewidth-3.5cm-5mm-2.9cm\relax}
    {\color{cursfosc}\small\bfseries\MakeUppercase{#3}}\par\vspace{2pt}
    {\color{navy}\fontsize{20}{23}\selectfont\bfseries #4\par}
  \end{minipage}\hfill
  \begin{minipage}[c]{2.6cm}\raggedleft\includegraphics[width=2.6cm]{\fm@logo}\end{minipage}\par
  \vspace{2pt}}

% ---------- barres de secció del color del curs ----------
\renewcommand{\seccio}[2]{\par\needspace{\fmneed\baselineskip}\bigskip
  \begin{tcolorbox}[colback=curs,colframe=curs,coltext=curstext,arc=2mm,boxrule=0pt,left=3mm,right=3mm,top=1.2mm,bottom=1.2mm]
  {\large\bfseries #1\quad #2}
  \end{tcolorbox}\vspace{-2pt}}

% ---------- estrelles de dificultat del color del curs (a Batxillerat no se'n posen) ----------
\@ifundefined{estrelles}{}{%
  \renewcommand{\estrelles}[1]{{\color{curs}\ifcase#1\or\faStar\faStar[regular]\faStar[regular]\or\faStar\faStar\faStar[regular]\or\faStar\faStar\faStar\fi}}}

% ---------- capçalera i peu de les pàgines ----------
% primera pàgina: logos institucionals al peu
\fancypagestyle{primera}{\fancyhf{}\renewcommand{\headrulewidth}{0pt}
  \fancyfoot[L]{\raisebox{-0.32\height}{\includegraphics[height=0.78cm]{logo-gva.png}}\hspace{5mm}%
    \raisebox{-0.32\height}{\includegraphics[height=0.78cm]{logo-vermellar.png}}}
  \fancyfoot[R]{\scriptsize\color{muted}\href{https://creativecommons.org/licenses/by-nc-sa/4.0/deed.ca}{\ccbyncsa}\ \ CC BY-NC-SA 4.0 · fentmates.org\qquad\small\thepage}}
% la resta: el símbol a l'esquerra i el curs a la dreta, amb el títol al centre com abans
\fancyhead[L]{\raisebox{-0.3\height}{\includegraphics[height=0.62cm]{\fm@marca}}}
\fancyhead[R]{\small\bfseries\color{cursfosc}\fm@cursnom}
\renewcommand{\headrule}{\color{curs!45}\hrule height \headrulewidth}
\makeatother

%   \fmcurs{1eso}
% i, al començament del document:
%   \capcalera{Matemàtiques}{1r ESO}{Fitxa d'exercicis · Unitat 2}{Divisibilitat}
% Opcional abans de \input: \def\fmneed{14} (espai mínim abans de cada secció).
% =====================================================================
\makeatletter
\def\fm@logodir{../../logo/pdf}
\providecommand{\fmneed}{10}
\def\fm@cursnom{}

% clau del curs, color i color del text damunt del color
\newcommand{\fm@curs}[3]{\@namedef{fm@color@#1}{#2}\@namedef{fm@text@#1}{#3}}
\fm@curs{1eso}{E6B84A}{navy}
\fm@curs{2eso}{3B82F6}{white}
\fm@curs{3eso}{FF7A50}{white}
\fm@curs{4eso-a}{A78BFA}{white}
\fm@curs{4eso-b}{6D28D9}{white}
\fm@curs{1batx-mi}{9B1C1C}{white}
\fm@curs{1batx-ccss}{EF6B6B}{white}
\fm@curs{2batx-mii}{1E6B47}{white}
\fm@curs{2batx-ccss}{2BB3A3}{white}

\newcommand{\fmcurs}[1]{%
  \@ifundefined{fm@color@#1}{\PackageError{fentmates}{Curs desconegut: #1}{}}{}%
  \edef\fm@tmp{\noexpand\definecolor{curs}{HTML}{\@nameuse{fm@color@#1}}%
               \noexpand\colorlet{curstext}{\@nameuse{fm@text@#1}}}\fm@tmp
  \colorlet{cursfosc}{curs!72!black}%
  \def\fm@logo{\fm@logodir/vertical-#1.pdf}%
  \def\fm@marca{\fm@logodir/marca-#1.pdf}}

% ---------- capçalera de la primera pàgina ----------
% #1 matèria (petit), #2 curs, #3 tipus de document i unitat, #4 títol
\newcommand{\capcalera}[4]{%
  \gdef\fm@cursnom{#2}%
  \par\noindent
  \begin{tcolorbox}[colback=curs,colframe=curs,boxrule=0pt,arc=2.2mm,width=3.5cm,height=2.15cm,
      valign=center,left=3mm,right=2mm,top=1mm,bottom=1mm,nobeforeafter,box align=center]
    \color{curstext}{\footnotesize #1}\par\vspace{1pt}{\fontsize{19}{21}\selectfont\bfseries #2\par}%
  \end{tcolorbox}\hspace{5mm}%
  \begin{minipage}[c]{\dimexpr\linewidth-3.5cm-5mm-2.9cm\relax}
    {\color{cursfosc}\small\bfseries\MakeUppercase{#3}}\par\vspace{2pt}
    {\color{navy}\fontsize{20}{23}\selectfont\bfseries #4\par}
  \end{minipage}\hfill
  \begin{minipage}[c]{2.6cm}\raggedleft\includegraphics[width=2.6cm]{\fm@logo}\end{minipage}\par
  \vspace{2pt}}

% ---------- barres de secció del color del curs ----------
\renewcommand{\seccio}[2]{\par\needspace{\fmneed\baselineskip}\bigskip
  \begin{tcolorbox}[colback=curs,colframe=curs,coltext=curstext,arc=2mm,boxrule=0pt,left=3mm,right=3mm,top=1.2mm,bottom=1.2mm]
  {\large\bfseries #1\quad #2}
  \end{tcolorbox}\vspace{-2pt}}

% ---------- estrelles de dificultat del color del curs (a Batxillerat no se'n posen) ----------
\@ifundefined{estrelles}{}{%
  \renewcommand{\estrelles}[1]{{\color{curs}\ifcase#1\or\faStar\faStar[regular]\faStar[regular]\or\faStar\faStar\faStar[regular]\or\faStar\faStar\faStar\fi}}}

% ---------- capçalera i peu de les pàgines ----------
% primera pàgina: logos institucionals al peu
\fancypagestyle{primera}{\fancyhf{}\renewcommand{\headrulewidth}{0pt}
  \fancyfoot[L]{\raisebox{-0.32\height}{\includegraphics[height=0.78cm]{logo-gva.png}}\hspace{5mm}%
    \raisebox{-0.32\height}{\includegraphics[height=0.78cm]{logo-vermellar.png}}}
  \fancyfoot[R]{\scriptsize\color{muted}\href{https://creativecommons.org/licenses/by-nc-sa/4.0/deed.ca}{\ccbyncsa}\ \ CC BY-NC-SA 4.0 · fentmates.org\qquad\small\thepage}}
% la resta: el símbol a l'esquerra i el curs a la dreta, amb el títol al centre com abans
\fancyhead[L]{\raisebox{-0.3\height}{\includegraphics[height=0.62cm]{\fm@marca}}}
\fancyhead[R]{\small\bfseries\color{cursfosc}\fm@cursnom}
\renewcommand{\headrule}{\color{curs!45}\hrule height \headrulewidth}
\makeatother

\fmcurs{2batx-mii}

\begin{document}
\thispagestyle{primera}

\capcalera{Matemáticas II}{2\textsuperscript{o} Bach.}{Soluciones de la ficha · Unidad 2}{Sistemas de ecuaciones lineales}
\vspace{8pt}

\begin{tcolorbox}[colback=soft,colframe=soft,arc=2mm,left=3mm,right=3mm,top=1.5mm,bottom=1.5mm]
\small Intenta resolver cada ejercicio antes de mirar su solución. Los pasos de Gauss pueden ser distintos de los tuyos y ser igualmente correctos; lo que debe coincidir es la clasificación y la solución (en un SCI, la solución general se puede escribir de maneras distintas según la incógnita que tomes como parámetro). Puedes comprobarlo con el resolutor de sistemas de los apuntes interactivos.
\end{tcolorbox}

% =====================================================
\seccio{1}{Sistemas de ecuaciones lineales. Método de Gauss}

\ex{1} \label{ex:g4}\[A^*=\left(\begin{array}{rrrr|r}1 & 1 & 1 & 1 & 2\\2 & -1 & 1 & -1 & 5\\1 & 2 & -1 & 3 & -3\\3 & 1 & -2 & 1 & -2\end{array}\right)\xrightarrow[\substack{F_3'=F_3-F_1\\F_4'=F_4-3F_1}]{F_2'=F_2-2F_1}\left(\begin{array}{rrrr|r}1 & 1 & 1 & 1 & 2\\0 & -3 & -1 & -3 & 1\\0 & 1 & -2 & 2 & -5\\0 & -2 & -5 & -2 & -8\end{array}\right)\]\[\xrightarrow{F_2\leftrightarrow F_3}\left(\begin{array}{rrrr|r}1 & 1 & 1 & 1 & 2\\0 & 1 & -2 & 2 & -5\\0 & -3 & -1 & -3 & 1\\0 & -2 & -5 & -2 & -8\end{array}\right)\]\[\xrightarrow[F_4''=F_4'+2F_2']{F_3''=F_3'+3F_2'}\left(\begin{array}{rrrr|r}1 & 1 & 1 & 1 & 2\\0 & 1 & -2 & 2 & -5\\0 & 0 & -7 & 3 & -14\\0 & 0 & -9 & 2 & -18\end{array}\right)\]\[\xrightarrow{F_4'''=7F_4''-9F_3''}\left(\begin{array}{rrrr|r}1 & 1 & 1 & 1 & 2\\0 & 1 & -2 & 2 & -5\\0 & 0 & -7 & 3 & -14\\0 & 0 & 0 & -13 & 0\end{array}\right)\]Resolviendo de abajo arriba: \ $\res{x=1;\ y=-1;\ z=2;\ t=0}$

\ex{2} \begin{tasks}(1)
  \task \[A^*=\left(\begin{array}{rrrr|r}1 & 1 & -1 & 1 & 1\\2 & 1 & 1 & -1 & 3\\3 & 2 & 0 & 0 & 4\end{array}\right)\xrightarrow[F_3'=F_3-3F_1]{F_2'=F_2-2F_1}\left(\begin{array}{rrrr|r}1 & 1 & -1 & 1 & 1\\0 & -1 & 3 & -3 & 1\\0 & -1 & 3 & -3 & 1\end{array}\right)\]\[\xrightarrow{F_3''=F_3'-F_2'}\left(\begin{array}{rrrr|r}1 & 1 & -1 & 1 & 1\\0 & -1 & 3 & -3 & 1\\0 & 0 & 0 & 0 & 0\end{array}\right)\]Rango 2 con 4 incógnitas: \rest{SCI con dos parámetros}. Con $z=\lambda$, $t=\mu$: \ $y=-1+3\lambda-3\mu$, \ $x=1-y+z-t$:\par $\res{x=2-2\lambda+2\mu;\ y=-1+3\lambda-3\mu;\ z=\lambda;\ t=\mu}$
  \task \[A^*=\left(\begin{array}{rrr|r}1 & 1 & 1 & 6\\1 & -1 & 2 & 5\\2 & 1 & -1 & 1\\3 & 2 & 0 & 7\end{array}\right)\xrightarrow[\substack{F_3'=F_3-2F_1\\F_4'=F_4-3F_1}]{F_2'=F_2-F_1}\left(\begin{array}{rrr|r}1 & 1 & 1 & 6\\0 & -2 & 1 & -1\\0 & -1 & -3 & -11\\0 & -1 & -3 & -11\end{array}\right)\]\[\xrightarrow{F_2\leftrightarrow F_3}\left(\begin{array}{rrr|r}1 & 1 & 1 & 6\\0 & -1 & -3 & -11\\0 & -2 & 1 & -1\\0 & -1 & -3 & -11\end{array}\right)\]\[\xrightarrow[F_4''=F_4'-F_2']{F_3''=F_3'-2F_2'}\left(\begin{array}{rrr|r}1 & 1 & 1 & 6\\0 & -1 & -3 & -11\\0 & 0 & 7 & 21\\0 & 0 & 0 & 0\end{array}\right)\]La 4.ª fila se anula: quedan tres ecuaciones independientes. \rest{SCD}: \ $\res{x=1;\ y=2;\ z=3}$
  \task Los mismos pasos dan la última fila $(0\ 0\ 0\,|\,1)$: \ \rest{SI}. (Las tres primeras ecuaciones obligan a $x=1$, $y=2$, $z=3$, y entonces $3x+2y=7\neq8$.)
\end{tasks}

\ex{3} \[\left(\begin{array}{rrr|r}1 & 2 & -1 & 1\\ 2 & 1 & 1 & 4\\ 4 & 5 & k & m\end{array}\right)\xrightarrow[F_3'=F_3-4F_1]{F_2'=F_2-2F_1}\left(\begin{array}{rrr|r}1 & 2 & -1 & 1\\ 0 & -3 & 3 & 2\\ 0 & -3 & k+4 & m-4\end{array}\right)\xrightarrow{F_3''=F_3'-F_2'}\left(\begin{array}{rrr|r}1 & 2 & -1 & 1\\ 0 & -3 & 3 & 2\\ 0 & 0 & k+1 & m-6\end{array}\right)\]\begin{tasks}(1)
  \task Si $k\neq-1$: \ \rest{SCD} (para cualquier $m$).
  \task Si $k=-1$ y $m=6$: la 3.ª fila es nula: \ \rest{SCI}.
  \task Si $k=-1$ y $m\neq6$: la 3.ª fila dice $0=m-6\neq0$: \ \rest{SI}.
\end{tasks}Para $k=-1$, $m=6$, con $z=\lambda$: \ $-3y+3\lambda=2\Rightarrow y=\lambda-\tfrac23$; \ $x=1-2y+z$:\par $\res{x=\tfrac73-\lambda;\ y=-\tfrac23+\lambda;\ z=\lambda}$

\ex{2} Hay infinitas respuestas. Buscamos relaciones que cumplan todas las soluciones: $x+y=3$ e $y+z=2$. La tercera ecuación debe ser combinación lineal de las dos (si no, el sistema sería SCD); por ejemplo, la suma: $x+2y+z=5$:\[\boxed{\left\{\begin{aligned}x+y&=3\\y+z&=2\\x+2y+z&=5\end{aligned}\right.}\]

\ex{2} \pau{Andalucía, extraordinaria 2025}\[\left\{\begin{aligned}x+y+z&=22\\4x+3y+2z&=61\end{aligned}\right.\] \ Restando a la 2.ª el doble de la 1.ª: $2x+y=17$. Con $x=\lambda$: \ $y=17-2\lambda$, \ $z=5+\lambda$ (\rest{SCI}).\par Para que sean naturales, $\lambda\geq1$ y $17-2\lambda\geq1$: $\lambda=1,\dots,8$. El 15 solo puede ser $y$ (con $\lambda=1$): \ \rest{1, 15 y 6}. Las otras opciones son $(2,13,7)$, $(3,11,8)$, $(4,9,9)$, $(5,7,10)$, $(6,5,11)$, $(7,3,12)$ y $(8,1,13)$.


% =====================================================
\seccio{2}{Expresión matricial. Método matricial}

\ex{2} $|A|=5$, \ $\mathrm{Adj}(A)=\begin{pmatrix}2 & 2 & -1\\-4 & 1 & 2\\3 & -2 & 1\end{pmatrix}$, \ $A^{-1}=\frac15\begin{pmatrix}2 & -4 & 3\\2 & 1 & -2\\-1 & 2 & 1\end{pmatrix}$.\[X_1=\frac15\begin{pmatrix}2 & -4 & 3\\2 & 1 & -2\\-1 & 2 & 1\end{pmatrix}\begin{pmatrix}2\\2\\3\end{pmatrix}=\boxed{\begin{pmatrix}1\\ 0\\ 1\end{pmatrix}},\qquad X_2=\frac15\begin{pmatrix}2 & -4 & 3\\2 & 1 & -2\\-1 & 2 & 1\end{pmatrix}\begin{pmatrix}1\\1\\4\end{pmatrix}=\boxed{\begin{pmatrix}2\\ -1\\ 1\end{pmatrix}}\]

\ex{2} $AX-X=B\Rightarrow(A-I)X=B$, \ con $A-I=\begin{pmatrix}2 & 1 & 0\\1 & 2 & 1\\0 & 1 & 2\end{pmatrix}$, \ $|A-I|=4$, \ $(A-I)^{-1}=\frac14\begin{pmatrix}3 & -2 & 1\\-2 & 4 & -2\\1 & -2 & 3\end{pmatrix}$:\[X=\frac14\begin{pmatrix}3 & -2 & 1\\-2 & 4 & -2\\1 & -2 & 3\end{pmatrix}\begin{pmatrix}1\\0\\1\end{pmatrix}=\boxed{\begin{pmatrix}1\\ -1\\ 1\end{pmatrix}}\]

\ex{3} \begin{tasks}(1)
  \task $|A|=\begin{vmatrix}1 & 1 & 0\\1 & k & 1\\0 & 1 & 1\end{vmatrix}=k-1-1=k-2$: \ si \ $\res{k\neq2}$
  \task $\mathrm{Adj}(A)=\begin{pmatrix}k-1 & -1 & 1\\ -1 & 1 & -1\\ 1 & -1 & k-1\end{pmatrix}$ (es simétrica), \ $A^{-1}=\dfrac{1}{k-2}\begin{pmatrix}k-1 & -1 & 1\\ -1 & 1 & -1\\ 1 & -1 & k-1\end{pmatrix}$:\[X=\frac{1}{k-2}\begin{pmatrix}k-1+0+1\\ -1+0-1\\ 1+0+k-1\end{pmatrix}=\frac{1}{k-2}\begin{pmatrix}k\\ -2\\ k\end{pmatrix}\ \Rightarrow\ \boxed{x=z=\tfrac{k}{k-2};\ \ y=\tfrac{2}{2-k}}\]
  \task $|A|=0$: la suma de la 1.ª y la 3.ª ecuaciones da $x+2y+z=2$, que contradice la 2.ª ($x+2y+z=0$): \ \rest{SI}.
\end{tasks}

\ex{3} \begin{tasks}(1)
  \task Si $AX=B$, multiplicando por $A^{-1}$ por la izquierda, $X=A^{-1}B$: la solución existe y es única. Con $B=O$, $X=A^{-1}O=O$.
  \task $A(X_1-X_2)=AX_1-AX_2=B-B=O$; \ \ $A\big(X_1+t(X_2-X_1)\big)=B+t\,(B-B)=B$.
  \task Si tiene dos distintas, $X_1+t(X_2-X_1)$ da una solución distinta para cada $t$: tiene \rest{infinitas}.
\end{tasks}

\ex{2} \pau{La Rioja, extraordinaria 2024}$A^{-1}b$ es la solución $X$ del sistema $AX=b$ (cuando $A$ es invertible: $|A|=-a^2+a+2=-(a-2)(a+1)\neq0$, es decir, $a\neq2$ y $a\neq-1$). Observando el sistema, $x=y=z=1$ cumple las tres ecuaciones: $1-1+a=a$, $a+1-1=a$, $a+1+1=a+2$. Como la solución es única:\[A^{-1}b=\boxed{\begin{pmatrix}1\\ 1\\ 1\end{pmatrix}}\]


% =====================================================
\seccio{3}{Regla de Cramer}

\ex{2} $|A|=a^3-3a+2=(a-1)^2(a+2)$: es de Cramer si $a\neq1$ y $a\neq-2$. \ Sumando todas las columnas a la primera y sacando factor común:\[|A_x|=\begin{vmatrix}1 & 1 & 1\\ 1 & a & 1\\ 1 & 1 & a\end{vmatrix}=(a-1)^2,\qquad x=\frac{(a-1)^2}{(a-1)^2(a+2)}=\frac{1}{a+2}\]y, de la misma manera (o por la simetría del sistema), $|A_y|=|A_z|=(a-1)^2$: \ $\res{x=y=z=\tfrac{1}{a+2}}$

\ex{3} $|A|=-13$ \ y \ $|A_t|=\begin{vmatrix}1 & 1 & 1 & 2\\2 & -1 & 1 & 5\\1 & 2 & -1 & -3\\3 & 1 & -2 & -2\end{vmatrix}=0$: \ $t=\dfrac{0}{-13}=\res{0}$, \ como habíamos obtenido por Gauss.

\ex{2} La 3.ª ecuación es la suma de las dos primeras, y $\begin{vmatrix}1 & 1\\ 1 & -1\end{vmatrix}=-2\neq0$: \ $\left\{\begin{aligned}x+y&=2-z\\ x-y&=1-2z\end{aligned}\right.$\[x=\frac{\begin{vmatrix}2-z & 1\\ 1-2z & -1\end{vmatrix}}{-2}=\frac{3-3z}{2},\qquad y=\frac{\begin{vmatrix}1 & 2-z\\ 1 & 1-2z\end{vmatrix}}{-2}=\frac{1+z}{2}\]\[\boxed{x=\tfrac{3-3\lambda}{2};\ \ y=\tfrac{1+\lambda}{2};\ \ z=\lambda,\quad\lambda\in\mathbb{R}}\]

\ex{3} \pau{Islas Baleares, extraordinaria 2025}\begin{tasks}(1)
  \task $|A|=\begin{vmatrix}k & 1 & 0\\0 & k & 1\\3 & -1 & -1\end{vmatrix}=-k^2+k+3$, que se anula para $k=\frac{1\pm\sqrt{13}}{2}$. Si $k\neq\frac{1\pm\sqrt{13}}{2}$: \rest{SCD}. Si $k=\frac{1\pm\sqrt{13}}{2}$: $\begin{vmatrix}k & 1\\ 0 & k\end{vmatrix}=k^2\neq0$, así que $\operatorname{rg}(A)=2$, y el menor de $A^*$ \ $\begin{vmatrix}k & 1 & 1\\0 & k & 0\\3 & -1 & 0\end{vmatrix}=-3k\neq0$ da $\operatorname{rg}(A^*)=3$: \rest{SI}.
  \task Por Cramer, con $|A_x|=1-k$, $|A_y|=3$, $|A_z|=-3k$:\[\boxed{x=\frac{1-k}{-k^2+k+3};\quad y=\frac{3}{-k^2+k+3};\quad z=\frac{-3k}{-k^2+k+3}}\]
\end{tasks}

\ex{2} \pau{Cataluña, extraordinaria 2026}\begin{tasks}(1)
  \task $|A|=3k-3$. Si $k\neq1$: \rest{SCD}. Si $k=1$: $\begin{vmatrix}-1 & 2\\ 1 & 1\end{vmatrix}\neq0$, $\operatorname{rg}(A)=2$, y $3E_1=5E_3-E_2$, también en los términos independientes: \rest{SCI}.
  \task Si $k\neq1$, por Cramer: $|A_y|=\begin{vmatrix}2 & 3 & 1\\-1 & 1 & 2\\1 & 2 & 1\end{vmatrix}=0\Rightarrow y=0$, y entonces $-x+2z=1$, $x+z=2$: \ $\res{x=1;\ y=0;\ z=1}$, que no depende de $k$. Si $k=1$: restando la 1.ª y la 3.ª, $x=1$, e $y+z=1$: \ $\res{x=1;\ y=1-\lambda;\ z=\lambda}$
\end{tasks}


% =====================================================
\seccio{4}{Sistemas homogéneos}

\ex{2} $|A|=-a^3+3a-2=-(a-1)^2(a+2)$.\begin{tasks}(1)
  \task Si $a\neq1$ y $a\neq-2$: \ \rest{SCD}, solo la solución trivial \ $x=y=z=0$.
  \task Si $a=1$: las tres ecuaciones son $x+y+z=0$, rango 1: \ \rest{SCI con dos parámetros}:\par $\res{x=-\lambda-\mu;\ y=\lambda;\ z=\mu}$
  \task Si $a=-2$: rango 2 ($\begin{vmatrix}1 & 1\\ 1 & -2\end{vmatrix}\neq0$). De las dos primeras, $x+y=2z$ y $x-2y=-z$: $y=z$, $x=z$: \ $\res{x=y=z=\lambda}$
\end{tasks}

\ex{2} La 3.ª ecuación es la suma de las dos primeras y las dos primeras no son proporcionales: $\operatorname{rg}(A)=2$, con 4 incógnitas: \rest{dos parámetros}. De la 3.ª, $x=-\tfrac32z$; con $z=\lambda$, $t=\mu$, de la 1.ª, $y=-x-z-t$:\[\boxed{x=-\tfrac32\lambda;\ \ y=\tfrac12\lambda-\mu;\ \ z=\lambda;\ \ t=\mu}\]

\ex{2} Por ejemplo, $y=2x$ y $z=-x$: \ $\res{\left\{\begin{aligned}2x-y&=0\\x+z&=0\end{aligned}\right.}$ \ (el rango es 2, de modo que la solución depende de un parámetro). \ $k\lambda+2\lambda-3\lambda=0\Rightarrow\res{k=1}$

\ex{3} \begin{tasks}(1)
  \task \rest{Verdadera}: $\operatorname{rg}(A)\leq$ número de ecuaciones $<n$, y un sistema homogéneo siempre es compatible: es SCI.
  \task \rest{Verdadera}: si $|A|\neq0$, sería SCD y solo tendría la solución trivial.
  \task \rest{Verdadera}: $A(aX_1+bX_2)=aAX_1+bAX_2=aO+bO=O$.
\end{tasks}

\ex{2} \pau{La Rioja, ordinaria 2025}$|A|=2-4m$. Si $m\neq\tfrac12$: \rest{solo la solución trivial}. Para $\res{m=\tfrac12}$, $\operatorname{rg}(A)=2$ (\,$\begin{vmatrix}3 & 1\\ 1 & -1\end{vmatrix}\neq0$): de la 1.ª y la 3.ª, $4x=2z$ e $y=z-3x$:\[\boxed{x=\lambda;\ \ y=-\lambda;\ \ z=2\lambda,\quad\lambda\in\mathbb{R}}\]


% =====================================================
\seccio{5}{Teorema de Rouché-Frobenius}

\ex{2} \label{ex:r4}El menor de las tres primeras filas de $A$ vale $7\neq0$, y $\operatorname{rg}(A)=3$.\begin{tasks}(1)
  \task $|A^*|=0$: \ $\operatorname{rg}(A^*)=3=n$: \ \rest{SCD}. Resolviendo las tres primeras ecuaciones: \ $\res{x=y=z=1}$ \ (que también cumple la 4.ª).
  \task $|A^*|=7\neq0$: \ $\operatorname{rg}(A^*)=4\neq3$: \ \rest{SI}.
\end{tasks}

\ex{2} Las dos primeras columnas son proporcionales, y también la 3.ª y la 4.ª (con signo cambiado): $\operatorname{rg}(A)\leq2$, y $\begin{vmatrix}1 & -1\\ 2 & 1\end{vmatrix}=3\neq0$ da $\operatorname{rg}(A)=2$. La 3.ª ecuación es $E_2-E_1$, de modo que $\operatorname{rg}(A^*)=2<4$: \rest{SCI con dos parámetros}.\par Sumando las dos primeras: $3x+6y=6$; con $y=\lambda$, $t=\mu$: \ $\res{x=2-2\lambda;\ y=\lambda;\ z=-1+\mu;\ t=\mu}$

\ex{2} $|A|=\begin{vmatrix}1 & 1 & 1\\1 & -1 & 2\\3 & 1 & 4\end{vmatrix}=0$ y $\begin{vmatrix}1 & 1\\ 1 & -1\end{vmatrix}\neq0$: $\operatorname{rg}(A)=2$ siempre, de modo que \rest{nunca es SCD}. Como $F_3=2F_1+F_2$ en la matriz $A$, el sistema es compatible si y solo si los términos independientes cumplen la misma relación:\[\boxed{c=2a+b}\](y entonces es SCI; si no, $\operatorname{rg}(A^*)=3$ y es SI).

\ex{2} \begin{tasks}(1)
  \task \rest{Falsa}: puede ser $\operatorname{rg}(A^*)=4$ (es el caso del segundo sistema del ejercicio \ref{ex:r4}).
  \task \rest{Verdadera}: $\operatorname{rg}(A)\leq3<4=n$.
  \task \rest{Verdadera}: $\operatorname{rg}(A^*)=4$ y $\operatorname{rg}(A)\leq3$.
  \task \rest{Falsa}: puede ser SCI (si $\operatorname{rg}(A)=\operatorname{rg}(A^*)<3$).
\end{tasks}

\ex{2} \pau{Islas Baleares, extraordinaria 2026}$|A|=\begin{vmatrix}2 & 1 & 1\\4 & 5 & 7\\1 & 2 & 3\end{vmatrix}=0$ y $\begin{vmatrix}2 & 1\\ 4 & 5\end{vmatrix}=6\neq0$: $\operatorname{rg}(A)=2$. La 2.ª ecuación es la suma de la 1.ª y el doble de la 3.ª, también en los términos independientes ($6500+16000=22500$): $\operatorname{rg}(A^*)=2<3$. \ \rest{SCI}: \rest{no} se pueden determinar de manera única.


% =====================================================
\seccio{6}{Discusión de sistemas con parámetros}

\ex{2} $|A|=(a-1)^2(a+2)$.\begin{tasks}(1)
  \task Si $a\neq1$ y $a\neq-2$: \rest{SCD}. Por Cramer (o restando las ecuaciones dos a dos):\par $\res{x=-\tfrac{a+1}{a+2};\ \ y=\tfrac{1}{a+2};\ \ z=\tfrac{(a+1)^2}{a+2}}$
  \task Si $a=1$: las tres ecuaciones son $x+y+z=1$: \ \rest{SCI con dos parámetros}:\par $\res{x=1-\lambda-\mu;\ y=\lambda;\ z=\mu}$
  \task Si $a=-2$: sumando las tres ecuaciones, $0=1-2+4=3$: \ \rest{SI} (\,$\operatorname{rg}(A)=2$, $\operatorname{rg}(A^*)=3$).
\end{tasks}

\ex{2} $|A|=\begin{vmatrix}1 & 1 & 1\\1 & 2 & a\\2 & 3 & 4\end{vmatrix}=3-a$.\begin{tasks}(1)
  \task Si $a\neq3$: \ \rest{SCD}.
  \task Si $a=3$ y $b=3$: la 3.ª ecuación es la suma de las dos primeras: \rest{SCI}.
  \task Si $a=3$ y $b\neq3$: los coeficientes de la 3.ª son la suma de los de las otras, pero el término independiente no: \rest{SI}.
\end{tasks}Para $a=b=3$, restando las dos primeras: $y+2z=1$; con $z=\lambda$: \ $\res{x=\lambda;\ y=1-2\lambda;\ z=\lambda}$

\ex{2} $\operatorname{rg}(A)=2$ siempre ($\begin{vmatrix}1 & 2\\ 2 & -1\end{vmatrix}=-5\neq0$). \ $|A^*|=\begin{vmatrix}1 & 2 & 3\\2 & -1 & 1\\a & 1 & 2\end{vmatrix}=5a-5$.\begin{tasks}(1)
  \task Si $a\neq1$: $\operatorname{rg}(A^*)=3$: \ \rest{SI}.
  \task Si $a=1$: $\operatorname{rg}(A^*)=2=n$: \ \rest{SCD}: \ $\res{x=1;\ y=1}$
\end{tasks}

\ex{3} Sumando las dos primeras ecuaciones: $2x+2z=1+a$. Las dos primeras filas de $A$ no son proporcionales, así que $\operatorname{rg}(A)\geq2$.\begin{tasks}(1)
  \task Si $k\neq2$: el menor de las columnas 1.ª, 2.ª y 3.ª, $\begin{vmatrix}1 & 1 & 1\\1 & -1 & 1\\2 & 0 & k\end{vmatrix}=-k+2\neq0$: \ $\operatorname{rg}(A)=\operatorname{rg}(A^*)=3<4$: \ \rest{SCI con un parámetro}.
  \task Si $k=2$ y $a=2$: la 3.ª ecuación, $2x+2z=3$, coincide con la suma de las dos primeras: rango 2: \ \rest{SCI con dos parámetros}.
  \task Si $k=2$ y $a\neq2$: $2x+2z=3$ y $2x+2z=1+a$ se contradicen: \ \rest{SI}.
\end{tasks}

\ex{2} $|A|=\begin{vmatrix}1 & 1 & 1\\m & 1 & m-1\\1 & m & 1\end{vmatrix}=m-1$.\begin{tasks}(1)
  \task Si $m\neq1$: \ \rest{SCD}.
  \task Si $m=1$: la 1.ª ecuación es $x+y+z=2$ y la 3.ª, $x+y+z=1$: \ \rest{SI}.
  \task $m=0$: \ $x+y+z=1$, \ $y-z=0$, \ $x+z=1$ \ $\Rightarrow$ \ $\res{x=1;\ y=0;\ z=0}$
  \task $m=2$: \ $x+y+z=3$, \ $2x+y+z=2$, \ $x+2y+z=1$ \ $\Rightarrow$ \ $\res{x=-1;\ y=-2;\ z=6}$
\end{tasks}

\ex{3} $|A|=\begin{vmatrix}1 & 1 & 1\\1 & 2 & a\\1 & 4 & a^2\end{vmatrix}$ es un determinante de Vandermonde (traspuesto) con los valores $1$, $2$ y $a$: \ $|A|=(2-1)(a-1)(a-2)=(a-1)(a-2)$.\begin{tasks}(1)
  \task Si $a\neq1$ y $a\neq2$: \ \rest{SCD} para cualquier $b$. Por Cramer, los tres numeradores también son determinantes de Vandermonde:\[x=\frac{\begin{vmatrix}1 & 1 & 1\\b & 2 & a\\b^2 & 4 & a^2\end{vmatrix}}{(a-1)(a-2)}=\frac{(2-b)(a-b)(a-2)}{(a-1)(a-2)},\qquad y=\frac{\begin{vmatrix}1 & 1 & 1\\1 & b & a\\1 & b^2 & a^2\end{vmatrix}}{(a-1)(a-2)}=\frac{(b-1)(a-1)(a-b)}{(a-1)(a-2)}\]\[z=\frac{\begin{vmatrix}1 & 1 & 1\\1 & 2 & b\\1 & 4 & b^2\end{vmatrix}}{(a-1)(a-2)}\ \Rightarrow\ \boxed{x=\frac{(2-b)(a-b)}{a-1};\ \ y=\frac{(b-1)(a-b)}{a-2};\ \ z=\frac{(b-1)(b-2)}{(a-1)(a-2)}}\]
  \task Si $a=1$ o $a=2$: dos columnas de $A$ son iguales y $\operatorname{rg}(A)=2$ (\,$\begin{vmatrix}1 & 1\\ 1 & 2\end{vmatrix}\neq0$). El menor de $A^*$ formado por las columnas de $x$, de $y$ y de los términos independientes es \ $\begin{vmatrix}1 & 1 & 1\\1 & 2 & b\\1 & 4 & b^2\end{vmatrix}=(b-1)(b-2)$:\par si $b\neq1$ y $b\neq2$, \ $\operatorname{rg}(A^*)=3$: \ \rest{SI}; \ si $b=1$ o $b=2$, \ $\operatorname{rg}(A^*)=2$: \ \rest{SCI}.
  \task Soluciones de los cuatro casos SCI (con $\lambda\in\mathbb{R}$):\par $a=1,\ b=1$: \ $\res{x=1-\lambda;\ y=0;\ z=\lambda}$ \qquad $a=1,\ b=2$: \ $\res{x=-\lambda;\ y=1;\ z=\lambda}$\par $a=2,\ b=1$: \ $\res{x=1;\ y=-\lambda;\ z=\lambda}$ \qquad $a=2,\ b=2$: \ $\res{x=0;\ y=1-\lambda;\ z=\lambda}$
\end{tasks}

\ex{3} \[A^*\xrightarrow[F_3'=F_3-3F_1]{F_2'=F_2-2F_1}\left(\begin{array}{rrrr|r}1 & 2 & -1 & 1 & 2\\ 0 & -3 & 3 & -3 & -3\\ 0 & -3 & m^2-1 & -3 & m-5\end{array}\right)\xrightarrow{F_3''=F_3'-F_2'}\left(\begin{array}{rrrr|r}1 & 2 & -1 & 1 & 2\\ 0 & -3 & 3 & -3 & -3\\ 0 & 0 & m^2-4 & 0 & m-2\end{array}\right)\]\begin{tasks}(1)
  \task Si $m\neq\pm2$: \ $\operatorname{rg}(A)=\operatorname{rg}(A^*)=3<4$: \ \rest{SCI con un parámetro}. De la 3.ª fila, $z=\frac{m-2}{m^2-4}=\frac{1}{m+2}$; de la 2.ª, $y=1+z-t$; de la 1.ª, $x=2-2y+z-t=t-z$. Con $t=\lambda$:\[\boxed{x=\lambda-\tfrac{1}{m+2};\ \ y=1+\tfrac{1}{m+2}-\lambda;\ \ z=\tfrac{1}{m+2};\ \ t=\lambda}\]
  \task Si $m=2$: la 3.ª fila es nula: \ $\operatorname{rg}(A)=\operatorname{rg}(A^*)=2<4$: \ \rest{SCI con dos parámetros}. Con $z=\lambda$, $t=\mu$: \ $\res{x=\mu-\lambda;\ y=1+\lambda-\mu;\ z=\lambda;\ t=\mu}$
  \task Si $m=-2$: la 3.ª fila dice $0=-4$: \ $\operatorname{rg}(A)=2\neq\operatorname{rg}(A^*)=3$: \ \rest{SI}.
\end{tasks}

\ex{2} \pau{Andalucía, ordinaria 2026}\begin{tasks}(1)
  \task $|A|=3\big[(1-a)(3-a)-8\big]=3(a^2-4a-5)=3(a-5)(a+1)$. \ Si $a\neq5$ y $a\neq-1$: \rest{SCD}. \ Si $a=-1$: las dos últimas ecuaciones son $2y+2z=0$ y $4y+4z=-6$: \rest{SI}. \ Si $a=5$: son $-4y+2z=0$ y $4y-2z=0$ (proporcionales): \rest{SCI}.
  \task $3x-y=0$, \ $y+2z=0$, \ $4y+3z=-5$ \ $\Rightarrow$ \ $\res{x=-\tfrac23;\ y=-2;\ z=1}$
\end{tasks}

\ex{2} \pau{Aragón, extraordinaria 2025}\begin{tasks}(1)
  \task $|A|=m^3-16m^2-m+16=(m-16)(m^2-1)$. \ Si $m\neq16$ y $m\neq\pm1$: \rest{una solución}. \ Si $m=1$: la 3.ª ecuación es la 1.ª cambiada de signo: \rest{infinitas soluciones}. \ Si $m=-1$: la 1.ª y la 3.ª dicen $x+2y-z=-1$ y $x+2y-z=1$: \rest{ninguna solución}. \ Si $m=16$: $\operatorname{rg}(A)=2$ y el menor $\begin{vmatrix}1 & -1 & -1\\8 & -6 & -8\\-1 & 256 & 16\end{vmatrix}=30\neq0$ de $A^*$ da $\operatorname{rg}(A^*)=3$: \rest{ninguna solución}.
  \task $x+2y-z=-1$, \ $8x+y-6z=-8$. Restando a la 2.ª el séxtuple de la 1.ª: $2x-11y=-2$. Con $y=2\lambda$: $x=-1+11\lambda$, $z=x+2y+1=15\lambda$: \ $\res{x=-1+11\lambda;\ y=2\lambda;\ z=15\lambda}$
\end{tasks}

\ex{2} \pau{Canarias, ordinaria 2024}\begin{tasks}(1)
  \task $|A|=-k^2+2k+3=-(k-3)(k+1)$. \ Si $k\neq3$ y $k\neq-1$: \rest{SCD}. \ Si $k=3$: la 1.ª y la 3.ª ecuaciones tienen los mismos coeficientes ($3x+y-3z$) y términos independientes distintos: \rest{SI}. \ Si $k=-1$: lo mismo ocurre con la 2.ª y la 3.ª ($-x+y+z$): \rest{SI}.
  \task Restando la 1.ª y la 3.ª: $z=4$; entonces $4x+y=17$ y $-x+y=-8$: \ $\res{x=5;\ y=-3;\ z=4}$
\end{tasks}

\ex{2} \pau{Cantabria, ordinaria 2026}\begin{tasks}(1)
  \task $|A|=1-a^2$. \ Si $a\neq\pm1$: \rest{SCD}. \ Si $a=-1$: las dos últimas son $y-z=1$ y $-y+z=3$: \rest{SI}. \ Si $a=1$: las dos últimas son $y+z=1$: \rest{SCI}.
  \task $y+2z=1$, \ $2y+z=0$ \ $\Rightarrow$ \ $y=-\tfrac13$, $z=\tfrac23$: \ $\res{x=-\tfrac13;\ y=-\tfrac13;\ z=\tfrac23}$
  \task $y+z=1$ y $x=-(y+z)=-1$: \ $\res{x=-1;\ y=1-\lambda;\ z=\lambda}$
\end{tasks}

\ex{3} \pau{Cantabria, extraordinaria 2026}\begin{tasks}(1)
  \task Si $a\neq0$: el menor $\begin{vmatrix}1 & a\\ a^2 & 0\end{vmatrix}=-a^3\neq0$, y $\operatorname{rg}(A)=\operatorname{rg}(A^*)=2<3$: \rest{SCI}. Si $a=0$: queda $y=1$ y $0=0$: rango 1, \rest{SCI con dos parámetros}. Es decir, \rest{SCI para todo $a$}.
  \task $\res{x=\lambda;\ y=1;\ z=\mu}$, \ $\lambda,\mu\in\mathbb{R}$.
  \task $8x+y+2z=1$, \ $2x+4y=2$. Con $y=\lambda$: $x=1-2\lambda$, \ $z=\frac{1-8x-y}{2}$: \ $\res{x=1-2\lambda;\ y=\lambda;\ z=-\tfrac72+\tfrac{15}{2}\lambda}$
\end{tasks}

\ex{2} \pau{Castilla-La Mancha, extraordinaria 2025}\begin{tasks}(1)
  \task $|A|=2a-6$. \ Si $a\neq3$: \rest{SCD} (una solución). \ Si $a=3$: la 2.ª ecuación, $4x+2y=3$, es la suma de la 1.ª y la 3.ª: \rest{SCI} (infinitas soluciones).
  \task Sumando la 1.ª y la 3.ª: $x+2y=3$; con $4x+2y=0$: $x=-1$, $y=2$, $z=-1$: \ $\res{x=-1;\ y=2;\ z=-1}$
\end{tasks}

\ex{3} \pau{Castilla y León, extraordinaria 2026}$A-\alpha I=\begin{pmatrix}2-\alpha & 2 & 0\\ 0 & 2-\alpha & 0\\ 2 & 0 & 4-\alpha\end{pmatrix}$, \ $|A-\alpha I|=(2-\alpha)^2(4-\alpha)$.\begin{tasks}(1)
  \task Si $\alpha\neq2$ y $\alpha\neq4$: \ \rest{SCD}.
  \task Si $\alpha=2$: queda $2y=0$, $0=0$, $2x+2z=1$: rango 2: \ \rest{SCI} \ ($x=\lambda$, $y=0$, $z=\tfrac12-\lambda$).
  \task Si $\alpha=4$: la 3.ª columna de $A$ es nula y $\operatorname{rg}(A)=2$, pero $\begin{vmatrix}-2 & 2 & 0\\0 & -2 & 0\\2 & 0 & 1\end{vmatrix}=4\neq0$: $\operatorname{rg}(A^*)=3$: \ \rest{SI}.
\end{tasks}

\ex{2} \pau{Castilla y León, ordinaria 2026}\begin{tasks}(1)
  \task $|A|=m^2-7m+10=(m-2)(m-5)$. \ Si $m\neq2$ y $m\neq5$: \rest{SCD}. \ Si $m=5$: la 1.ª ecuación es la 3.ª cambiada de signo: \rest{SCI}. \ Si $m=2$: $\begin{vmatrix}-8 & -2\\ 7 & 4\end{vmatrix}\neq0$, $\operatorname{rg}(A)=2$, y $\begin{vmatrix}-8 & -2 & -1\\7 & 4 & 3\\5 & 2 & 1\end{vmatrix}=6\neq0$: $\operatorname{rg}(A^*)=3$, \rest{SI}.
  \task Sumando la 1.ª y la 3.ª: $-2x=0$, $x=0$; entonces $5y+2z=4$ y $2y+z=1$: \ $\res{x=0;\ y=2;\ z=-3}$
\end{tasks}

\ex{2} \pau{Cataluña, ordinaria 2026}\begin{tasks}(1)
  \task $|A|=m^2-4m+3=(m-1)(m-3)$. \ Si $m\neq1$ y $m\neq3$: \rest{SCD}. \ Si $m=1$: la 1.ª y la 3.ª ecuaciones coinciden: \rest{SCI}. \ Si $m=3$: la 1.ª y la 2.ª tienen los mismos coeficientes y términos independientes 4 y 5: \rest{SI}.
  \task $x+y+z=4$, $x+3y+z=5$: $y=\tfrac12$: \ $\res{x=\tfrac72-\lambda;\ y=\tfrac12;\ z=\lambda}$
  \task Sí: $\left(3,\tfrac12,\tfrac12\right)$ es la solución con $\lambda=\tfrac12$. Basta con añadir una ecuación que fije $\lambda$, por ejemplo \ $\res{z=\tfrac12}$: entonces $\operatorname{rg}(A)=\operatorname{rg}(A^*)=3$ y es SCD.
\end{tasks}

\ex{2} \pau{Extremadura, extraordinaria 2026}\begin{tasks}(1)
  \task $|A|=k^2-k-6=(k-3)(k+2)$. \ Si $k\neq3$ y $k\neq-2$: \rest{SCD}. \ Si $k=-2$: $\operatorname{rg}(A)=2$ y $\begin{vmatrix}-2 & -1 & -6\\2 & -2 & 6\\1 & 1 & -2\end{vmatrix}=-30\neq0$: \rest{SI}. \ Si $k=3$: la 1.ª ecuación es la suma de la 2.ª y la 3.ª: \rest{SCI}.
  \task $2x+3y-2z=6$, \ $x-y+z=3$. Sumando a la 1.ª el doble de la 2.ª: $4x+y=12$. Con $x=3-\lambda$: $y=4\lambda$, \ $z=3-x+y=5\lambda$: \ $\res{x=3-\lambda;\ y=4\lambda;\ z=5\lambda}$
\end{tasks}

\ex{2} \pau{Extremadura, extraordinaria 2025}\begin{tasks}(1)
  \task $|A|=m^2-m-2=(m-2)(m+1)$. \ Si $m\neq2$ y $m\neq-1$: \rest{SCD}. \ Si $m=2$: la 1.ª ecuación es el doble de la 2.ª más tres veces la 3.ª ($0=6-6$): \rest{SCI}. \ Si $m=-1$: $-E_2+6E_3$ tiene los mismos coeficientes que la 1.ª ecuación, $-x+7y+5z$, pero término independiente $-15\neq0$: \rest{SI}.
  \task Restando la 3.ª a la 2.ª: $x=5$; entonces $7y+5z=-5$ e $y+z=-2$: \ $\res{x=5;\ y=\tfrac52;\ z=-\tfrac92}$
\end{tasks}

\ex{2} \pau{Galicia, extraordinaria 2026}$|A|=-2m-4=-2(m+2)$. \ Si $m\neq-2$: \rest{SCD}. \ Si $m=-2$: restando a la 1.ª el doble de la 2.ª se obtiene $2y-2z=5$, mientras que la 3.ª es $2y-2z=-7$: \rest{SI}.

\ex{2} \pau{Galicia, extraordinaria 2025}$|A|=m^2-4m=m(m-4)$. \ Si $m\neq0$ y $m\neq4$: \rest{SCD}. \ Si $m=0$: queda $x+z=0$, $x+3z=0$, $2z=0$: $x=z=0$ e $y$ cualquiera: \rest{SCI}. \ Si $m=4$: restando las dos primeras, $4y+2z=-4$, mientras que la 3.ª es $4y+2z=12$: \rest{SI}.

\ex{2} \pau{Madrid, extraordinaria 2025}\begin{tasks}(1)
  \task El sistema es $AX=\begin{pmatrix}0\\ k\\ 2k\end{pmatrix}$, con $|A|=k^3+2k^2+k=k(k+1)^2$. \ Si $k\neq0$ y $k\neq-1$: \rest{SCD}. \ Si $k=0$: es homogéneo con $\operatorname{rg}(A)=2$: \rest{SCI}. \ Si $k=-1$: la 1.ª más la 3.ª da $y+z=-2$, mientras que la 2.ª es $y+z=-1$: \rest{SI}.
  \task $y+z=0$, \ $x+y=0$: \ $\res{x=-\lambda;\ y=\lambda;\ z=-\lambda}$
\end{tasks}

\ex{3} \pau{Navarra, extraordinaria 2026}Utilizamos el \textbf{teorema de Rouché-Frobenius}, que relaciona el tipo de sistema con los rangos de $A$ y de $A^*$. \ $|A|=(a^2-4)^2=(a-2)^2(a+2)^2$.\begin{tasks}(1)
  \task Si $a\neq\pm2$: $\operatorname{rg}(A)=\operatorname{rg}(A^*)=3$: \ \rest{SCD}.
  \task Si $a=2$: la 1.ª y la 3.ª ecuaciones son $4y+2z=4$ y $4y+2z=-4$: \rest{SI}.
  \task Si $a=-2$: las ecuaciones quedan $2z=4$, $z=2$, $-2z=-4$: $\operatorname{rg}(A)=\operatorname{rg}(A^*)=1<3$: \rest{SCI}: \ $\res{x=\lambda;\ y=\mu;\ z=2}$
\end{tasks}

\ex{3} \pau{Navarra, extraordinaria 2025}Teorema de Rouché-Frobenius. Sumando la 2.ª y la 3.ª: $(m+2)z=2m+4$. \ $|A|=m(m-1)^2(m+2)$.\begin{tasks}(1)
  \task Si $m\neq0,1,-2$: \rest{SCD}: $z=2$; restando la 1.ª a la 2.ª, $(m-1)y+z=1+m^2$, de donde $y=m+1$, y de la 2.ª se despeja $x$: \ $\res{x=\tfrac{1-3m-2m^2}{m(m-1)};\ y=m+1;\ z=2}$
  \task Si $m=0$: la columna de $x$ es nula, $\operatorname{rg}(A)=2$, y $\begin{vmatrix}2 & -1 & 1\\1 & 0 & 2\\-1 & 2 & 2\end{vmatrix}=-2\neq0$: \rest{SI}.
  \task Si $m=1$: $3y-z=0$, $3y=2$, $-3y+3z=4$: \rest{SCI}: \ $\res{x=\lambda;\ y=\tfrac23;\ z=2}$
  \task Si $m=-2$: la 3.ª ecuación es la 2.ª cambiada de signo: \rest{SCI}: de $6x-z=-3$ y $6x-3y=2$: \ $\res{x=\lambda;\ y=-\tfrac23+2\lambda;\ z=3+6\lambda}$
\end{tasks}

\ex{2} \pau{País Vasco, extraordinaria 2025}\begin{tasks}(1)
  \task Ordenando: $\left\{\begin{aligned}x-\alpha y+3z&=3\\ \alpha x+y+z&=2\\ x-y+2z&=1\end{aligned}\right.$. \ $|A|=2\alpha^2-4\alpha=2\alpha(\alpha-2)$. \ Si $\alpha\neq0$ y $\alpha\neq2$: \rest{SCD}. \ Si $\alpha=0$: la 1.ª ecuación es la suma de las otras dos: \rest{SCI}. \ Si $\alpha=2$: \rest{SI} ($\operatorname{rg}(A)=2$, $\operatorname{rg}(A^*)=3$).
  \task $y+z=2$, \ $x-y+2z=1$: con $z=\lambda$: \ $\res{x=3-3\lambda;\ y=2-\lambda;\ z=\lambda}$
\end{tasks}


% =====================================================
\seccio{7}{Problemas de planteamiento}

\ex{2} $(100z+10y+x)-(100x+10y+z)=396\Leftrightarrow-99x+99z=396$:\[\left\{\begin{aligned}x+y+z&=12\\-99x+99z&=396\\x-y+z&=0\end{aligned}\right.\]De la 1.ª y la 3.ª: $2y=12$, $y=6$; entonces $x+z=6$ y $z-x=4$: \ el número es \rest{165}.

\ex{2} Sustituyendo cada punto: \ \[\left\{\begin{aligned}-a+b-c+d&=-4\\d&=1\\a+b+c+d&=2\\8a+4b+2c+d&=11\end{aligned}\right.\] \ $d=1$; sumando la 1.ª y la 3.ª: $2b+2d=-2$, $b=-2$; entonces $a+c=3$ y $8a+2c=18$: $a=2$, $c=1$.\[\boxed{p(x)=2x^3-2x^2+x+1}\]

\ex{2} \begin{tasks}(1)
  \task $x$, $y$, $z$: monedas de 1, 2 y 5 €. \ \[\left\{\begin{aligned}x+y+z&=10\\x+2y+5z&=25\end{aligned}\right.\] \ Rango 2 y 3 incógnitas: \rest{SCI}. Restando: $y+4z=15$; con $z=\lambda$: \ $\res{x=-5+3\lambda;\ y=15-4\lambda;\ z=\lambda}$
  \task Deben ser naturales: $x\geq0\Rightarrow\lambda\geq\tfrac53$ e $y\geq0\Rightarrow\lambda\leq\tfrac{15}{4}$, de modo que $\lambda=2$ o $\lambda=3$: \ \rest{1, 7 y 2} \ o \ \rest{4, 3 y 3} monedas.
  \task \rest{4 monedas de 1 €, 3 de 2 € y 3 de 5 €}.
\end{tasks}

\ex{3} \begin{tasks}(1)
  \task $x$ generales, $y$ infantiles, $z$ sénior: \ \[\left\{\begin{aligned}x+y+z&=100\\10x+5y+kz&=600\\y-z&=0\end{aligned}\right.\] \ $|A|=\begin{vmatrix}1 & 1 & 1\\10 & 5 & k\\0 & 1 & -1\end{vmatrix}=k-15$. Si $k\neq15$: \rest{SCD}; si $k=15$: con $y=z$, $x+2y=100$ y $10x+20y=600$ se contradicen: \rest{SI}.
  \task \rest{20 generales, 40 infantiles y 40 sénior}.
  \task Con $k=10$ el sistema es SCD, pero la solución es $x=-60$, $y=z=80$: \rest{no tiene sentido} (no puede haber un número negativo de entradas).
\end{tasks}

\ex{2} \pau{Aragón, extraordinaria 2024}$x$, $y$, $z$: gramos de M$_1$, M$_2$ y M$_3$. Pasando a mg y dividiendo entre 10:\[\left\{\begin{aligned}x+y+2z&=7000\\x+2y+3z&=9000\\2x+3y+5z&=16000\end{aligned}\right.\]La 3.ª ecuación es la suma de las otras dos: \rest{SCI}. Restando las dos primeras, $y+z=2000$; con $z=\lambda$: \ $\res{x=5000-\lambda;\ y=2000-\lambda;\ z=\lambda}$, \ con \rest{$0\leq\lambda\leq2000$} para que las cantidades no sean negativas. Es posible, y de muchas maneras.

\ex{2} \pau{Asturias, extraordinaria 2026}\begin{tasks}(1)
  \task \[\left\{\begin{aligned}x+y+z&=10\\x+2y+3z&=m+9\\2x+3y+(7-m)z&=22\end{aligned}\right.\]
  \task $|A|=3-m$: hace falta \ $\res{m=3}$. Entonces la 3.ª ecuación es la suma de las otras dos: \rest{SCI}:\par $\res{x=8+\lambda;\ y=2-2\lambda;\ z=\lambda}$
  \task $\res{x=11;\ y=-2;\ z=1}$: \rest{no tiene sentido}: el número de camiones no puede ser negativo.
\end{tasks}

\ex{3} \pau{Asturias, ordinaria 2026}\begin{tasks}(1)
  \task $c$, $t$, $s$: precios del café, el té y el zumo: \ \[\left\{\begin{aligned}c+2t+2s&=11\\c+3t+3s&=16\\3c+3t+ps&=21\end{aligned}\right.\] \ $|A|=p-3$. Si $p\neq3$: \rest{SCD}; si $p=3$: \rest{SI} (restando las dos primeras, $t+s=5$, y entonces la 3.ª daría $3c=6$ pero la 1.ª, $c=1$).
  \task Con $p=2$: \ $c=1$, $t=8$, $s=-3$: un precio negativo: \rest{tiene razón el cliente}.
  \task Con $p=4$: \ \rest{café 1 €, té 2 € y zumo 3 €}.
\end{tasks}

\ex{2} \pau{Canarias, extraordinaria 2025}$x$, $y$, $z$: kg de cada materia prima. Hay cuatro ecuaciones (masa total, proteína, grasa y carbohidratos):\[\left\{\begin{aligned}x+y+z&=1000\\0{,}2x+0{,}4y+0{,}6z&=360\\0{,}1x+0{,}2y+0{,}1z&=120\\0{,}1x+0{,}3y+0{,}3z&=200\end{aligned}\right.\]$|A^*|=0$ y el menor de las tres primeras filas de $A$ es no nulo: $\operatorname{rg}(A)=\operatorname{rg}(A^*)=3$, \rest{SCD}. Resolviendo: \ \rest{500 kg de vegetales, 200 kg de harinas y 300 kg de cárnicos}.

\ex{2} \pau{Castilla-La Mancha, extraordinaria 2024}\begin{tasks}(1)
  \task \[\left\{\begin{aligned}x+y+z&=157\\x+2y+3z&=278\\x-kz&=0\end{aligned}\right.\]
  \task $|A|=1-k$: solución única si \ $\res{k\neq1}$ \ ($k>0$). Vender el mismo número quiere decir $k=1$, y entonces el sistema es incompatible ($x=z$ da $2x+y=157$ y $4x+2y=278$, contradictorias): \rest{no es posible}.
\end{tasks}

\ex{2} \pau{Madrid, ordinaria 2025}$x$, $y$, $z$: lanzamientos de 1, 2 y 3 puntos: \ $0{,}8x+0{,}5\cdot2y+0{,}4\cdot3z=80$, \ $\frac y3=\frac{x+z}{5}$, \ $2z=x+y-5$:\[\left\{\begin{aligned}0{,}8x+y+1{,}2z&=80\\3x-5y+3z&=0\\x+y-2z&=5\end{aligned}\right.\]$\Rightarrow$ \ \rest{25 lanzamientos de 1 punto, 30 de 2 puntos y 25 de 3 puntos}.

\ex{2} \pau{Murcia, extraordinaria 2026}\begin{tasks}(1)
  \task \[\left\{\begin{aligned}x+2y+z&=400\\x+y+2z&=500\\x+y+z&=k\end{aligned}\right.\]
  \task $|A|=1\neq0$: \rest{SCD para todo $k$}. Restando la 3.ª a las otras: $y=400-k$, $z=500-k$, y $x=k-y-z$: \ $\res{x=3k-900;\ y=400-k;\ z=500-k}$. \ Para que tenga sentido, $k$ debe ser un número natural entre 300 y 400.
\end{tasks}

\ex{2} \pau{Murcia, ordinaria 2026}\begin{tasks}(1)
  \task \[\left\{\begin{aligned}x+y+z&=12\\x+2y-z&=k\\2x-y+z&=4\end{aligned}\right.\]
  \task $|A|=-7\neq0$ para todo $k$: \rest{SCD}. Por Cramer: \ $\res{x=\tfrac{2k}{7};\ y=4+\tfrac{k}{7};\ z=8-\tfrac{3k}{7}}$
\end{tasks}

\ex{2} \pau{País Vasco, extraordinaria 2026}\begin{tasks}(1)
  \task $a$, $b$, $c$: recipientes de cada productor: \ \[\left\{\begin{aligned}10a+15b+12c&=612\\a+b+c&=48\\5a+9b+4c&=304\end{aligned}\right.\]
  \task Solución: $a=12$, $b=20$, $c=16$: \ \rest{60 € a A, 180 € a B y 64 € a C}.
\end{tasks}


\end{document}
